\documentclass[11pt]{article}
\usepackage[margin=1in]{geometry}
\usepackage{amsmath}
\usepackage{amssymb}
\usepackage{amsthm}
\usepackage{enumitem}
\usepackage{xcolor}
\usepackage{tcolorbox}

\theoremstyle{definition}
\newtheorem{definition}{Definition}
\newtheorem{theorem}{Theorem}

\title{ECO 310: Problem 2a(ii) Detailed Explanation\\Risk Aversion and Certainty Equivalents}
\author{Explanation Notes}
\date{February 28, 2026}

\begin{document}

\maketitle

\section{Problem Statement}

\begin{tcolorbox}[colback=blue!5!white,colframe=blue!75!black,title=Problem 2a(ii)]
If Mavi is more risk-averse than Melsi, then from Lecture 10, her risk premium is higher for every gamble. \textbf{Prove that this implies that for a particular gamble, Melsi's certainty equivalent is higher than Mavi's.}
\end{tcolorbox}

\section{Key Definitions}

Before proving the result, let's establish the fundamental concepts:

\begin{definition}[Certainty Equivalent (CE)]
The certainty equivalent of a gamble is the guaranteed amount of wealth that provides the same utility as the expected utility of the gamble:
\[
\text{EU}(gamble) = u(\text{CE})
\]
\end{definition}

\begin{definition}[Risk Premium ($\pi$)]
The risk premium is the maximum amount a decision-maker would pay to avoid the risk and receive the expected value with certainty:
\[
\text{EU}(gamble) = u(E[W] - \pi)
\]
where $E[W]$ is the expected wealth from the gamble.
\end{definition}

\begin{definition}[Risk Aversion]
A person is risk-averse if they prefer a certain outcome to a gamble with the same expected value. More risk-averse individuals have higher risk premiums for any given gamble.
\end{definition}

\section{The Fundamental Relationship}

\subsection{Step 1: Connecting CE and $\pi$}

We have two equations that both equal EU(gamble):
\begin{align}
\text{EU}(gamble) &= u(\text{CE}) \label{eq:ce}\\
\text{EU}(gamble) &= u(E[W] - \pi) \label{eq:rp}
\end{align}

Since both expressions equal EU(gamble), we can equate them:
\[
u(\text{CE}) = u(E[W] - \pi)
\]

\subsection{Step 2: Using Monotonicity of Utility}

Since utility functions are strictly increasing in wealth (more money is always better), the function $u(\cdot)$ is one-to-one. This means:
\[
u(\text{CE}) = u(E[W] - \pi) \implies \text{CE} = E[W] - \pi
\]

\subsection{Step 3: Rearranging to Isolate CE}

From the equation $\text{CE} = E[W] - \pi$, we can rearrange:
\begin{equation}
\boxed{\text{CE} = E[W] - \pi} \label{eq:key}
\end{equation}

This is our \textbf{key relationship}: the certainty equivalent equals the expected value minus the risk premium.

\section{The Proof}

Now we can prove the main result.

\subsection{Setup}

Consider a particular gamble with expected wealth $E[W]$.

Let:
\begin{itemize}
\item $\text{CE}_{\text{Melsi}}$ = Melsi's certainty equivalent for this gamble
\item $\text{CE}_{\text{Mavi}}$ = Mavi's certainty equivalent for this gamble
\item $\pi_{\text{Melsi}}$ = Melsi's risk premium for this gamble
\item $\pi_{\text{Mavi}}$ = Mavi's risk premium for this gamble
\end{itemize}

\subsection{Given Information}

We are told that Mavi is more risk-averse than Melsi, which means:
\begin{equation}
\pi_{\text{Mavi}} > \pi_{\text{Melsi}} \label{eq:given}
\end{equation}

\subsection{What We Need to Prove}

We want to show:
\[
\text{CE}_{\text{Melsi}} > \text{CE}_{\text{Mavi}}
\]

\subsection{Step-by-Step Proof}

\begin{proof}

\textbf{Step 1:} Apply equation \eqref{eq:key} to Melsi:
\[
\text{CE}_{\text{Melsi}} = E[W] - \pi_{\text{Melsi}}
\]

\textbf{Step 2:} Apply equation \eqref{eq:key} to Mavi:
\[
\text{CE}_{\text{Mavi}} = E[W] - \pi_{\text{Mavi}}
\]

\textbf{Step 3:} Note that both dogs face the same gamble, so the expected wealth $E[W]$ is identical for both.

\textbf{Step 4:} Take the difference between Melsi's and Mavi's certainty equivalents:
\begin{align*}
\text{CE}_{\text{Melsi}} - \text{CE}_{\text{Mavi}} &= (E[W] - \pi_{\text{Melsi}}) - (E[W] - \pi_{\text{Mavi}})\\
&= E[W] - \pi_{\text{Melsi}} - E[W] + \pi_{\text{Mavi}}\\
&= \pi_{\text{Mavi}} - \pi_{\text{Melsi}}
\end{align*}

\textbf{Step 5:} From our given information \eqref{eq:given}, we know $\pi_{\text{Mavi}} > \pi_{\text{Melsi}}$, which means:
\[
\pi_{\text{Mavi}} - \pi_{\text{Melsi}} > 0
\]

\textbf{Step 6:} Therefore:
\[
\text{CE}_{\text{Melsi}} - \text{CE}_{\text{Mavi}} > 0
\]

\textbf{Conclusion:}
\[
\boxed{\text{CE}_{\text{Melsi}} > \text{CE}_{\text{Mavi}}}
\]

\end{proof}

\section{Intuition}

Why does this make sense?

\begin{itemize}[leftmargin=*]
\item The \textbf{certainty equivalent} is the guaranteed amount you'd accept instead of taking the gamble.

\item The \textbf{risk premium} is how much you'd pay (or how much you're discounting the expected value) to avoid the risk.

\item A more risk-averse person (Mavi) demands a bigger discount for bearing risk, so her risk premium $\pi$ is higher.

\item Since $\text{CE} = E[W] - \pi$, a higher risk premium means a lower certainty equivalent.

\item In other words: \textbf{Mavi needs more compensation to accept the same level of risk}, so the guaranteed amount that would make her indifferent is lower.

\item Conversely, Melsi (less risk-averse) is more willing to accept risk, so she'd be satisfied with a higher certain payment.
\end{itemize}

\section{Connection to Part (i)}

This result connects to part (i) of the problem, which showed that someone accepts a bet if and only if $\text{CE} > w_0$ (initial wealth).

Since:
\begin{itemize}
\item Both dogs are equally wealthy (same $w_0$)
\item Whenever Mavi accepts a bet, her $\text{CE}_{\text{Mavi}} > w_0$
\item Melsi's CE is even higher: $\text{CE}_{\text{Melsi}} > \text{CE}_{\text{Mavi}} > w_0$
\item Therefore Melsi also accepts the bet
\end{itemize}

\textbf{Conclusion:} Any bet that Mavi (more risk-averse) is willing to accept, Melsi (less risk-averse) will also accept.

\end{document}
